\addcontentsline{toc}{chapter}{Abstract}

\chapter*{Abstract}

\vspace*{-8mm}

Physics-Informed Neural Networks with sinusoidal activation functions
(SIREN-PINNs) are evaluated for solving the Helmholtz equation in one and two
dimensions as a foundation for optical speckle simulation. With $\omega_0 = 1.0$
calibrated to the wavenumber, the SIREN architecture achieves relative $L^2$
errors of $0.006\%$ in 1D ($5 \times 64$ neurons) and $0.171\%$ in 2D
($5 \times 128$) against analytical solutions, three and nearly two orders of
magnitude below the $5\%$ thesis threshold.
The speckle extension uses a spatially correlated random-phase screen and a
semianalytical angular-spectrum reference, and is solved through a modal
reduction that imposes the Cauchy condition exactly within the architecture.
Across five independent phase realizations at $z=1\lambda$ the complex-field
$L^2$ error is $3.161\%$, dominated by the evanescent fraction of the spectrum
that the modal basis cannot represent; measured over the propagating subspace
the error is $0.042\%$. The ensemble contrast is $C=0.9489$, against a mean
reference ensemble contrast of $C=0.9661$ over 64 realizations each.
A two-phase Adam~+~L-BFGS optimization strategy and Latin Hypercube Sampling
(LHS) are used in the training protocol.
Reproducibility was verified with three independent random seeds, yielding a
component-wise mean $L^2$ error of $0.192 \pm 0.089\%$ (population standard
deviation).

\bigskip

\noindent\textbf{Keywords:} Physics-Informed Neural Networks, SIREN,
Helmholtz equation, optical speckle, sinusoidal activation, Latin Hypercube
Sampling.
